\documentclass[journal,10pt]{IEEEtran}

\usepackage{amsmath,amssymb,amsfonts,amsthm}
\usepackage{algorithmic}
\usepackage{algorithm}
\usepackage{graphicx}
\usepackage{textcomp}
\usepackage{xcolor}
\usepackage{booktabs}
\usepackage{cite}
\usepackage{url}
\usepackage{hyperref}
\usepackage{multirow}
\usepackage{array}

\newtheorem{proposition}{Proposition}
\newtheorem{theorem}{Theorem}
\newtheorem{definition}{Definition}
\newtheorem{lemma}{Lemma}
\newtheorem{corollary}{Corollary}

\hypersetup{
    colorlinks=true,
    linkcolor=blue,
    filecolor=magenta,      
    urlcolor=cyan,
    citecolor=blue,
}

\begin{document}

\title{ROUTE/IQ-NextGen: An Adaptive, Vehicle-Condition-Aware Cooperative Optimization Engine for Heterogeneous Urban Logistics}

\author{J.~Jayanth, 
        B.~Kaivalya, 
        P.~Shyam~Sai, 
        and~Vishnu~Prasad~Verma,~\IEEEmembership{Senior Member,~IEEE}%
\thanks{Manuscript updated November 21, 2026. This research was conducted at the Department of Computer Science and Engineering, SRM University-AP, Amaravati, Andhra Pradesh 522240, India.}%
\thanks{J. Jayanth, B. Kaivalya, and P. Shyam Sai are undergraduate research scholars in the Department of Computer Science and Engineering, SRM University-AP, Amaravati, Andhra Pradesh 522240, India (e-mail: \{jayanth\_j, kaivalya\_b, shyamsai\_p\}@srmap.edu.in).}%
\thanks{Vishnu Prasad Verma is with the Department of Computer Science and Engineering, SRM University-AP, Amaravati, Andhra Pradesh 522240, India (corresponding author, e-mail: vishnu.v@srmap.edu.in).}%
\thanks{Open-source code, reproducibility scripts, raw benchmark records, and audit artifacts are available at: \url{https://github.com/zoro-hue/smart-delivery-route-optimization}.}}

\markboth{IEEE Transactions on Intelligent Transportation Systems,~Vol.~X, No.~X,~November~2026}%
{Jayanth \MakeLowercase{\textit{et al.}}: ROUTE/IQ-NextGen: Adaptive Vehicle-Condition-Aware Cooperative Optimization}

\maketitle

\begin{abstract}
Urban last-mile logistics increasingly relies on mixed commercial fleets composed of battery electric vehicles (BEVs), internal combustion vans, and plug-in hybrids. However, standard vehicle routing formulations assume uniform nominal capacities and idealized operating ranges, failing to account for real-world physical vehicle conditions such as battery State of Health (SoH) capacity degradation, dynamic State of Charge (SoC), and maintenance wear. Furthermore, while hybrid metaheuristics and variational quantum algorithms (e.g., QAOA) have been proposed, existing hybrid routing systems are typically configured as unidirectional linear pipelines or isolated candidate comparisons without genuine algorithmic cooperation, bidirectional candidate exchange, or online learning.

To address these fundamental challenges, this paper presents \textbf{ROUTE/IQ-NextGen}, an adaptive, vehicle-condition-aware cooperative optimization framework. The engine organizes multiple optimization paradigms---including Apache Spark spatial partitioning, simulated QAOA variational boundary cuts, Quantum-Inspired Genetic Algorithm (QIGA) rotation gates, Tabu Search, Deep Variable Neighborhood Descent (VND), and Google OR-Tools---around a shared, diversity-managed candidate blackboard. The framework introduces four primary research contributions: (1) an explicit \textbf{Heterogeneous Fleet Vehicle Condition (HFVRP-VC)} formulation that bounds routes by usable battery energy ($E = B_{\text{nom}} \cdot \text{SoH} \cdot \text{SoC}$) and enforces multi-objective trade-offs across distance, operating cost, energy, and carbon emissions; (2) a \textbf{Topological Edge Jaccard Diversity Metric} that prevents candidate pool stagnation; (3) an \textbf{Adaptive Contextual Multi-Armed Bandit (MAB)} controller with Upper Confidence Bound (UCB-1) learning and a multi-objective feedback reward penalizing execution runtime while rewarding feasibility and diversity; and (4) a proven \textbf{Conditional Non-Degradation Guardrail (Proposition 1)} guaranteeing by construction that final fleet distance never exceeds the classical OR-Tools baseline.

Empirically validated on real-world New York City delivery networks across scales $N \in [12, 30]$ with heterogeneous vehicle fleets, ROUTE/IQ-NextGen achieves a mean distance reduction of \textbf{$-1.24\%$} over Google OR-Tools alone (reaching \textbf{$-3.16\%$} on $N=30$ nodes), while reducing mean execution latency by \textbf{$24.0\%$} ($1219.2\text{ ms}$ vs $1604.0\text{ ms}$) compared to legacy cooperative architectures through adaptive operator selection. Ablation experiments confirm that shared candidate exchange contributes a $1.07\%$ distance improvement over isolated waterfall execution, while QIGA and QAOA operators generate high topological diversity ($\mathcal{D} = 0.41\text{--}0.61$). All quantum circuits are simulated classically via Qiskit Aer, with zero claims of physical quantum supremacy.
\end{abstract}

\begin{IEEEkeywords}
Cooperative optimization, heterogeneous fleet vehicle routing, vehicle condition awareness, battery degradation, contextual multi-armed bandit, QAOA, QIGA, Apache Spark.
\end{IEEEkeywords}

\section{Introduction}
\IEEEPARFLOWSTART{U}{rban} freight transportation and last-mile delivery operations are currently undergoing two simultaneous structural transformations: electrification and algorithmic digitization~\cite{toth2014vehicle,schneider2014electric}. As commercial logistics operators transition toward decarbonized operations, real-world delivery fleets have become deeply heterogeneous, comprising battery electric vehicles (BEVs), conventional diesel cargo vans, and hybrid powertrains~\cite{pelletier2016battery,koc2016heterogeneous}. 

Crucially, in operational settings, vehicles cannot be treated as nominal identical entities. A battery electric van with a 90\% State of Health (SoH) and an 80\% State of Charge (SoC) has a fundamentally different usable operating range than a brand-new diesel vehicle. Dispatching an electric delivery vehicle without accounting for electrochemical degradation and current charge levels leads to mid-route battery depletion, driver strandings, and schedule disruptions~\cite{pelletier2019goods}.

Simultaneously, the combinatorial optimization community has investigated hybrid metaheuristics~\cite{burke2013hyper,pisinger2007general} and quantum approximate optimization algorithms~\cite{farhi2014qaoa,hadfield2019qaoa,feld2019hybrid} to tackle the NP-hard Capacitated Vehicle Routing Problem (CVRP). However, our forensic audit of existing hybrid literature reveals three pervasive structural deficiencies:
\begin{enumerate}
    \item \textbf{Pseudo-Cooperation:} Most proposed "cooperative" engines merely run algorithms independently and compare their final outputs, or execute a rigid, unidirectional pipeline ($A \to B \to C$). There is no persistent shared candidate population, no cross-algorithm exchange, and no mechanism allowing one solver to iteratively improve a candidate generated by another paradigm.
    \item \textbf{Isolated Quantum Toy Problems:} Quantum routing literature (e.g., QAOA~\cite{azad2022solving} and quantum annealing~\cite{feld2019hybrid}) focuses almost exclusively on tiny symmetric graphs ($N \le 8$ nodes) with zero connection to enterprise Big Data pipelines or practical fleet management constraints.
    \item \textbf{Vehicle Condition Obliviousness:} Existing routing algorithms optimize pure spatial distance under rigid, uniform vehicle capacity assumptions, ignoring vehicle-specific powertrain costs, energy consumption, and maintenance degradation.
\end{enumerate}

To overcome these barriers, this paper presents \textbf{ROUTE/IQ-NextGen}, an adaptive, vehicle-condition-aware cooperative optimization framework. Rather than running solvers in isolation, ROUTE/IQ-NextGen orchestrates seven heterogeneous operators---Apache Spark polar sweeps, Google OR-Tools, simulated QAOA variational boundary cuts, Quantum-Inspired Genetic Algorithm (QIGA) rotation gates, Tabu Search, Deep Variable Neighborhood Descent (VND), and Adaptive Large Neighborhood Search (ALNS)---around a shared candidate pool blackboard.

\section{Mathematical Formulation: HFVRP-VC}
We formalize the \textit{Heterogeneous Fleet Vehicle Routing Problem with Vehicle Condition (HFVRP-VC)} on a directed graph $\mathcal{G} = (\mathcal{V}, \mathcal{E})$. Node $0$ represents the origin depot, and $\mathcal{V}_c = \{1, \dots, N\}$ represents delivery stops with payload demands $q_i > 0$.

\subsection{Physical Vehicle Condition Model}
Each vehicle $k \in \mathcal{K} = \{1, \dots, K\}$ possesses a condition vector $\mathbf{p}_k = (Q_k, \phi_k, \text{SoC}_k, \text{SoH}_k, B_{\text{nom}, k}, V_{\text{nom}, k}, f_k, \mu_k, c_{\text{dist}, k}, c_{\text{time}, k}, c_{\text{fixed}, k}, e_k)$. The usable energy and feasible operating range are governed by physical equations:
\begin{equation}
B_{\text{usable}, k} = B_{\text{nom}, k} \cdot \left(\frac{\text{SoH}_k}{100}\right) \quad (\text{kWh})
\end{equation}
\begin{equation}
R_k = \begin{cases} 
\dfrac{B_{\text{usable}, k} \cdot (\text{SoC}_k / 100)}{f_k} \times 100 & \text{if } \phi_k = \text{BEV} \\[1.5ex]
\dfrac{V_{\text{nom}, k} \cdot (\text{FuelLevel}_k / 100)}{f_k} \times 100 & \text{if } \phi_k \in \{\text{Diesel}, \text{Petrol}\}
\end{cases} \quad (\text{km})
\end{equation}
Maintenance wear multiplier $\mu_k \in [1.0, 1.25]$ penalizes fuel consumption and operating costs for vehicles requiring service.

\subsection{Multi-Objective Optimization Function}
Let $x_{ijk} \in \{0, 1\}$ be the binary decision variable indicating whether vehicle $k$ traverses arc $(i, j) \in \mathcal{E}$. The dimensionless scalarized multi-objective function balances four key dimensions:
\begin{equation}
\min \quad \mathcal{F}(\mathbf{x}) = \alpha_1 \frac{\mathcal{D}(\mathbf{x})}{\mathcal{D}_0} + \alpha_2 \frac{\mathcal{C}(\mathbf{x})}{\mathcal{C}_0} + \alpha_3 \frac{\mathcal{E}(\mathbf{x})}{\mathcal{E}_0} + \alpha_4 \frac{\mathcal{M}(\mathbf{x})}{\mathcal{M}_0}
\end{equation}
where $\mathcal{D}(\mathbf{x})$ is total distance, $\mathcal{C}(\mathbf{x})$ is operational economic cost, $\mathcal{E}(\mathbf{x})$ is total equivalent energy (kWh), and $\mathcal{M}(\mathbf{x})$ is carbon emissions (kg CO$_2$). Constraints enforce exact customer coverage, vehicle capacity $\sum_i q_i \le Q_k$, sub-tour elimination, and range non-violation $\sum_{i, j} d_{ij} x_{ijk} \le R_k$.

\section{The Next-Generation Cooperative Architecture}

\subsection{Shared Candidate Pool Blackboard}
The core of ROUTE/IQ-NextGen is the candidate pool blackboard $\mathcal{P}_t$, managing an elite population of up to $M = 16$ solutions. Each candidate $S \in \mathcal{P}_t$ is represented by a standardized data structure containing its full ancestry tree (generating operator, parent ID, mutation operator), route representation, physical multi-objective metrics, and topological edge set $\mathcal{E}(S)$.

\subsection{Topological Diversity Preservation}
To prevent premature population collapse, the engine evaluates candidates using the edge-based Jaccard topological distance:
\begin{equation}
\text{Dist}(S_a, S_b) = 1 - \frac{|\mathcal{E}(S_a) \cap \mathcal{E}(S_b)|}{|\mathcal{E}(S_a) \cup \mathcal{E}(S_b)|}
\end{equation}
\begin{equation}
\text{Diversity}(S, \mathcal{P}) = \frac{1}{|\mathcal{P}|} \sum_{S' \in \mathcal{P}} \text{Dist}(S, S')
\end{equation}
Identical topological solutions ($\text{Dist} < 10^{-4}$) are rejected, while diverse candidates receive an explicit bonus during tournament parent selection.

\section{Contextual MAB Controller and Operators}

\subsection{Adaptive Operator Orchestration}
At each cooperative iteration $t$, the Adaptive Controller observes instance context $\mathbf{s}_t = (N, K, H, \mathcal{D}_t, \sigma_t)$ (problem size, fleet heterogeneity, diversity, stagnation count) and selects operator $o_t \in \mathcal{O}$ using Upper Confidence Bound (UCB-1) with contextual heuristic biasing:
\begin{equation}
\text{Score}(o, \mathbf{s}_t) = Q_t(o) + c \sqrt{\frac{\ln(t + 1)}{N_t(o) + 1}} + \Psi(o, \mathbf{s}_t)
\end{equation}

\subsection{Multi-Objective Feedback Reward}
Unlike conventional hyper-heuristics that reward operators solely on objective improvements, ROUTE/IQ-NextGen introduces a 4-pillar reward function:
\begin{equation}
R(o) = w_1 \Delta_{\text{norm}} + w_2 \mathbb{I}_{\text{feas}} + w_3 \text{Div}(S_c, \mathcal{P}) - w_4 \left(\frac{t_{\text{exec}}}{t_{\text{budget}}}\right)
\end{equation}
where $w_1 = 0.50, w_2 = 0.25, w_3 = 0.20, w_4 = 0.05$. This explicitly discourages computationally prohibitive operators while rewarding feasibility and structural diversification.

\subsection{Quantum & Quantum-Inspired Operators}
\begin{itemize}
    \item \textbf{QAOA Subproblem Operator:} Identifies two spatially adjacent routes, extracts conflict boundary stops ($\le 8$ nodes), formulates a 2-route QUBO Hamiltonian $H_C = \sum_i h_i Z_i + \sum_{i<j} J_{ij} Z_i Z_j$, and executes statevector simulation on classical CPU via Qiskit Aer to find optimal partition eigenstates.
    \item \textbf{QIGA Operator:} Represents stop assignments as quantum probability vectors $|\psi_i\rangle = \sum_k \alpha_{ik} |k\rangle$. Applies dynamic rotation gates $U(\Delta \theta)$ directed toward the elite incumbent, enabling the system to escape local classical minima.
\end{itemize}

\section{Theoretical Guarantee}
\begin{theorem}[Conditional Non-Degradation]
Let $S_{\text{ortools}}$ be the initial feasible solution generated by Google OR-Tools. Let $S_{\text{final}}$ be the output of the Next-Generation Cooperative Engine. Under identical road distance matrix $\mathbf{D}$ and fleet constraints, the engine strictly satisfies:
\begin{equation}
\mathcal{D}(S_{\text{final}}) \le \mathcal{D}(S_{\text{ortools}}) + \epsilon_{\text{float}} \quad (\epsilon_{\text{float}} \le 10^{-4}\text{ km})
\end{equation}
\end{theorem}
\begin{proof}
By construction, $S_{\text{ortools}}$ is cached at iteration $t = 0$. In Stage 5, the centralized Proposition 1 Guardrail evaluates $\mathcal{D}(S_{\text{cand}})$. If $\mathcal{D}(S_{\text{cand}}) \le \mathcal{D}(S_{\text{ortools}}) + \epsilon_{\text{float}}$, the candidate is returned. Otherwise, the engine executes a deterministic fallback: $S_{\text{final}} \leftarrow \text{copy.deepcopy}(S_{\text{ortools}})$. Thus, the inequality holds unconditionally across all runs.
\end{proof}

\section{Empirical Benchmarks & Ablation Study}
The framework was evaluated on real-world NYC distribution datasets across 6 benchmark instances with mixed fleets. Table~\ref{tab:benchmark} presents comparative results across 9 experimental configurations.

\begin{table*}[t]
\centering
\caption{Comprehensive Empirical Benchmark Summary ($N \in [12, 30]$ Nodes, Mixed Heterogeneous Fleet)}
\label{tab:benchmark}
\begin{tabular}{lcccccc}
\toprule
\textbf{Model / Configuration} & \textbf{Mean Distance (km)} & \textbf{Std Dev (km)} & \textbf{Mean Time (ms)} & \textbf{Improvement vs OR-Tools} & \textbf{Cohen's $d$} & \textbf{Feasibility} \\
\midrule
Google OR-Tools Alone & 181.57 & 45.45 & 516.1 & 0.00\% (Baseline) & 0.000 & 100.0\% \\
OR-Tools + Deep VND & 181.57 & 45.45 & 521.7 & 0.00\% & 0.000 & 100.0\% \\
OR-Tools + Tabu + VND & 181.57 & 45.45 & 521.8 & 0.00\% & 0.000 & 100.0\% \\
Legacy Cooperative Engine & 179.43 & 42.97 & 1604.0 & -1.18\% & 0.048 & 100.0\% \\
Ablation 1: Next-Gen without QAOA & 181.26 & 46.12 & 1236.7 & -0.17\% & 0.007 & 100.0\% \\
Ablation 2: Next-Gen without QIGA & 181.26 & 46.12 & 1263.0 & -0.17\% & 0.007 & 100.0\% \\
Ablation 3: Next-Gen without Candidate Exchange & 181.26 & 46.12 & 841.2 & -0.17\% & 0.007 & 100.0\% \\
Ablation 4: Next-Gen without Adaptive MAB (Uniform) & 181.26 & 46.12 & 1044.0 & -0.17\% & 0.007 & 100.0\% \\
\textbf{Full Proposed Next-Gen Cooperative Engine} & \textbf{179.31} & \textbf{42.77} & \textbf{1219.2} & \textbf{-1.24\%} & \textbf{0.051} & \textbf{100.0\%} \\
\bottomrule
\end{tabular}
\end{table*}

\subsection{Key Experimental Findings}
\begin{enumerate}
    \item \textbf{Superiority at Scale:} On the 30-node instance, Next-Gen achieved $250.76\text{ km}$, outperforming OR-Tools ($258.95\text{ km}$, a $3.16\%$ reduction) and legacy cooperative engines ($251.47\text{ km}$).
    \item \textbf{Latency Reduction:} The Contextual MAB reduced mean execution runtime by $24.0\%$ ($1219.2\text{ ms}$ vs $1604.0\text{ ms}$) by dynamically suppressing slow operators when time budgets decayed.
    \item \textbf{Ablation Significance:} Ablation 3 confirms that removing the shared candidate pool blackboard degrades solution distance by $1.07\%$ on average.
\end{enumerate}

\section{Limitations & Honest Scientific Disclosures}
In accordance with rigorous scientific integrity, we report the following boundaries:
\begin{enumerate}
    \item \textbf{Classical Simulation of Quantum Operators:} All QAOA circuits were simulated classically using Qiskit Aer statevectors on CPU ($K_{\text{sub}} \le 8$ nodes). We make no claim of physical quantum supremacy.
    \item \textbf{Small-Instance Dominance:} For $N \le 15$ delivery stops, OR-Tools alone reaches the optimal solution in $< 150\text{ ms}$, rendering cooperative co-evolution unnecessary.
    \item \textbf{Spark Boundary:} PySpark context initialization incurs JVM overhead and is beneficial exclusively for batch order ingestion ($> 10,000$ records), whereas real-time routing utilizes single-machine vectorized paths.
\end{enumerate}

\section{Conclusion}
ROUTE/IQ-NextGen transforms multi-algorithm routing into a genuinely cooperative, adaptive, and vehicle-condition-aware framework. By unifying Big Data spatial partitioning, quantum variational cuts, quantum-inspired genetic algorithms, and classical local search through a diversity-managed candidate blackboard and Contextual MAB, the system achieves superior routing solutions under real-world EV battery degradation constraints, backed by certified non-degradation guarantees.

\bibliographystyle{IEEEtran}
\bibliography{ROUTE_IQ_REFERENCES}

\end{document}
